\chapter{エントロピー正則化}
\label{ch:seminar-03-entropic}

%% ============================================================
%% 3.1 エントロピー正則化
%% ============================================================
\section{エントロピー正則化}
\label{sec:sem-entropic-regularization}

\begin{definition}{離散エントロピー}{sem-discrete-entropy}
  非負行列 $\mathbf{P} \in \Rnn^{n \times m}$ に対して，
  \textbf{離散エントロピー}を
  \[
    \Hb(\mathbf{P})
    \defeq
    - \sum_{i=1}^{n} \sum_{j=1}^{m}
      P_{i,j}\bigl(\log P_{i,j} - 1\bigr)
  \]
  で定める．ただし $0\log 0 = 0$ と約束する．
\end{definition}

\begin{remark}{通常の Shannon エントロピーとの違い}{sem-entropy-convention}
  確率行列では $\sum_{i,j}P_{i,j}=1$ なので，
  \[
    \Hb(\mathbf{P})
    =
    -\sum_{i,j} P_{i,j}\log P_{i,j} + 1
  \]
  であり，通常の Shannon エントロピーとは定数 $1$ だけ異なる．
  この定数は最適化の解には影響しない．この形を使う理由は，
  $-\Hb$ の微分が
  \[
    \frac{\partial(-\Hb)}{\partial P_{i,j}} = \log P_{i,j}
  \]
  と簡単になるからである．
\end{remark}

\begin{remark}{$\varphi(x)=x\log x-x$ の連続性}{sem-phi-continuous}
  $-\Hb$ の被加数 $\varphi(x) \defeq x\log x - x$ は $[0,\infty)$ 上で連続である．
  実際，$(0,\infty)$ 上では連続関数の積・和として連続であり，
  $x = 0$ では，約束 $0\log 0 = 0$ により $\varphi(0) = 0$ と定めると，
  $\lim_{x\to0+} x\log x = 0$ から
  \[
    \lim_{x \to 0+} \varphi(x)
    = \lim_{x \to 0+}\bigl(x\log x - x\bigr) = 0 = \varphi(0)
  \]
  となって $x = 0$ でも（右）連続である．
  ゆえに $-\Hb(\mathbf{P}) = \sum_{i,j}\varphi(P_{i,j})$ は
  $\Rnn^{n\times m}$ 上で連続である．
\end{remark}

\begin{definition}{エントロピー正則化された離散最適輸送}{sem-entropic-ot}
  $\mathbf{a} \in \Rpos^n$（$\sum_{i=1}^n a_i = 1$），
  $\mathbf{b} \in \Rpos^m$（$\sum_{j=1}^m b_j = 1$），
  $\mathbf{C} \in \Rnn^{n \times m}$，$\varepsilon > 0$ とし，
  \[
    \CouplingsD(\mathbf{a}, \mathbf{b})
    \defeq
    \left\{
      \mathbf{P} \in \Rnn^{n \times m}
      \;\middle|\;
      \mathbf{P}\ones_m = \mathbf{a},\;
      \mathbf{P}^\top\ones_n = \mathbf{b}
    \right\}
  \]
  とおく．\textbf{エントロピー正則化された離散 Kantorovich 問題}を
  \[
    \MKD_{\mathbf{C}}^\varepsilon(\mathbf{a}, \mathbf{b})
    \defeq
    \min_{\mathbf{P} \in \CouplingsD(\mathbf{a}, \mathbf{b})}
    \left\{
      \inner{\mathbf{C}}{\mathbf{P}}
      - \varepsilon \Hb(\mathbf{P})
    \right\}
  \]
  と定める．
\end{definition}

\begin{claim}{負エントロピーの狭義凸性}{sem-neg-entropy-strict-convex}
  負エントロピー
  \[
    -\Hb(\mathbf{P}) = \sum_{i,j}\varphi(P_{i,j}),
    \qquad \varphi(x) \defeq x\log x - x
  \]
  は $\CouplingsD(\mathbf{a}, \mathbf{b})$ 上で狭義凸である．
\end{claim}

\begin{proof}
  $-\Hb(\mathbf{P}) = \sum_{i,j}\varphi(P_{i,j})$（$\varphi(x) = x\log x - x$）であるから，
  \textbf{(1)} $\varphi$ が $[0,\infty)$ 上で狭義凸であること，
  \textbf{(2)} そこから和 $\sum_{i,j}\varphi(P_{i,j})$ が狭義凸になること，
  の順に示す．

  \textbf{(1) $\varphi$ の狭義凸性．}
  $\varphi$ は $(0,\infty)$ 上で $C^2$ 級かつ $\varphi''(x) = 1/x > 0$ をみたすから，
  2階導関数の符号による判定で（厳密に）狭義凸である．
  ただし以下ではこの判定を用いず，導関数 $\varphi' = \log$ が狭義単調増加であることと
  平均値定理のみで示す．こちらは $\varphi$ が $C^1$ 級で導関数が狭義単調増加でありさえ
  すれば（$\varphi''$ の存在＝$C^2$ 級を要さず）成り立つため，仮定が弱く一般性が高い．
  $\varphi$ は $[0,\infty)$ 上で連続（注意~\ref{rem:sem-phi-continuous}），
  $(0,\infty)$ 上で微分可能で，
  $\varphi'(x) = \log x$ は $(0,\infty)$ 上で狭義単調増加である．
  相異なる $x_1, x_2 \in [0,\infty)$ と $\lambda \in (0,1)$ をとり，
  $x^* \defeq \lambda x_1 + (1-\lambda)x_2 \in (0,\infty)$ とおく．
  目標は
  \[
    \varphi(x^*) < \lambda\varphi(x_1) + (1-\lambda)\varphi(x_2)
  \]
  を示すことである．

  $\varphi$ は閉区間 $[x_1, x^*]$, $[x^*, x_2]$ 上で連続，その内部で微分可能だから，
  平均値定理（定理~\ref{thm:sem-mvt}）より
  \[
    \varphi(x_1) - \varphi(x^*) = \varphi'(c_1)\,(x_1 - x^*),
    \qquad
    \varphi(x_2) - \varphi(x^*) = \varphi'(c_2)\,(x_2 - x^*)
  \]
  をみたす $c_1$（$x_1$ と $x^*$ の間），$c_2$（$x^*$ と $x_2$ の間）が存在する．
  $x_1 = 0$ のときも $c_1 \in (0, x^*)$ は内点なので $\varphi'(c_1) = \log c_1$ は定義される．

  $x^* = \lambda x_1 + (1-\lambda)x_2$ より
  $x_1 - x^* = (1-\lambda)(x_1 - x_2)$,
  $x_2 - x^* = -\lambda(x_1 - x_2)$
  である．これらを上の 2 式に代入し，第 1 式を $\lambda$ 倍，第 2 式を $(1-\lambda)$ 倍して加えると
  \begin{align*}
    \lambda\varphi(x_1) + (1-\lambda)\varphi(x_2) - \varphi(x^*)
    &= \lambda(1-\lambda)(x_1 - x_2)\,\varphi'(c_1)
       - (1-\lambda)\lambda(x_1 - x_2)\,\varphi'(c_2) \\
    &= \lambda(1-\lambda)(x_1 - x_2)\bigl(\varphi'(c_1) - \varphi'(c_2)\bigr)
  \end{align*}
  を得る（左辺は $\lambda + (1-\lambda) = 1$ より $\varphi(x^*)$ がまとまる）．

  最後にこの符号を調べる．$c_1, c_2$ は $x^*$ をはさんで反対側にあるから
  $c_1 - c_2$ は $x_1 - x_2$ と同符号であり，$\varphi'$ が狭義単調増加だから
  $\varphi'(c_1) - \varphi'(c_2)$ も $c_1 - c_2$（ひいては $x_1 - x_2$）と同符号である．
  ゆえに $(x_1 - x_2)\bigl(\varphi'(c_1) - \varphi'(c_2)\bigr) > 0$ となり，
  $\lambda(1-\lambda) > 0$ と合わせて右辺は正．したがって目標の不等式が成り立ち，
  $\varphi$ は $[0,\infty)$ 上で狭義凸である．

  \textbf{(2) $-\Hb$ の狭義凸性．}
  相異なる $\mathbf{P}, \mathbf{Q} \in \CouplingsD(\mathbf{a}, \mathbf{b})$ と $t \in (0,1)$ をとり，
  凸結合 $(1-t)\mathbf{P} + t\mathbf{Q}$ における $-\Hb$ を評価する．
  $-\Hb(\mathbf{R}) = \sum_{i,j}\varphi(R_{i,j})$ であり，
  行列 $(1-t)\mathbf{P} + t\mathbf{Q}$ の $(i,j)$ 成分は $(1-t)P_{i,j} + tQ_{i,j}$ だから，
  \begin{align*}
    -\Hb\bigl((1-t)\mathbf{P} + t\mathbf{Q}\bigr)
    &= \sum_{i,j}\varphi\bigl((1-t)P_{i,j} + tQ_{i,j}\bigr) \\
    &\leq \sum_{i,j}\bigl[(1-t)\varphi(P_{i,j}) + t\varphi(Q_{i,j})\bigr] \\
    &= (1-t)\sum_{i,j}\varphi(P_{i,j}) + t\sum_{i,j}\varphi(Q_{i,j}) \\
    &= (1-t)\bigl(-\Hb(\mathbf{P})\bigr) + t\bigl(-\Hb(\mathbf{Q})\bigr)
  \end{align*}
  となる．等号の各行は $-\Hb = \sum_{i,j}\varphi$ の定義と和の線形性（スカラーのくくり出し）による．
  第 2 行の不等号は，各成分に (1) の狭義凸性
  $\varphi\bigl((1-t)P_{i,j}+tQ_{i,j}\bigr) \leq (1-t)\varphi(P_{i,j}) + t\varphi(Q_{i,j})$
  を適用したもので，$P_{i,j} \neq Q_{i,j}$ の成分では狭義（$<$），
  $P_{i,j} = Q_{i,j}$ の成分では等号である．
  $\mathbf{P} \neq \mathbf{Q}$ より狭義となる成分が少なくとも 1 つあるので，
  上の不等号は全体として狭義であり，
  \[
    -\Hb\bigl((1-t)\mathbf{P} + t\mathbf{Q}\bigr)
    < (1-t)\bigl(-\Hb(\mathbf{P})\bigr) + t\bigl(-\Hb(\mathbf{Q})\bigr).
  \]
  よって $-\Hb$ は $\CouplingsD(\mathbf{a}, \mathbf{b})$ 上で狭義凸である．
\end{proof}

\begin{proposition}{正則化解の存在と一意性}{sem-entropic-existence-unique}
  任意の $\varepsilon > 0$ に対して，
  $\MKD_{\mathbf{C}}^\varepsilon(\mathbf{a}, \mathbf{b})$ は一意な最適解
  \[
    \mathbf{P}_\varepsilon
    \defeq
    \argmin_{\mathbf{P} \in \CouplingsD(\mathbf{a}, \mathbf{b})}
    \left\{
      \inner{\mathbf{C}}{\mathbf{P}}
      - \varepsilon \Hb(\mathbf{P})
    \right\}
    \;\in\; \CouplingsD(\mathbf{a}, \mathbf{b})
  \]
  を持つ．
\end{proposition}

\begin{proof}
  \textbf{存在．}
  $\CouplingsD(\mathbf{a}, \mathbf{b})$ は空でないコンパクト集合である
  （主張~\ref{clm:sem-discrete-kantorovich-existence} の証明）．
  $\varphi(x) = x\log x - x$ は $[0,\infty)$ 上で連続（注意~\ref{rem:sem-phi-continuous}）だから，
  目的関数
  \[
    F(\mathbf{P})
    \defeq
    \inner{\mathbf{C}}{\mathbf{P}} - \varepsilon \Hb(\mathbf{P})
    =
    \inner{\mathbf{C}}{\mathbf{P}}
    + \varepsilon \sum_{i,j} \varphi(P_{i,j})
  \]
  も連続である．
  コンパクト集合上の連続関数は最小値を達成するから，
  最適解が存在する．

  \textbf{一意性．}
  正則化項がコストの線形性を破り，目的関数 $F$ を狭義凸化する点が鍵である．
  $\mathbf{P}^\star, \mathbf{Q}^\star \in \CouplingsD(\mathbf{a}, \mathbf{b})$ をともに最適解とし，
  最適値を $m \defeq \MKD_{\mathbf{C}}^\varepsilon(\mathbf{a}, \mathbf{b}) = F(\mathbf{P}^\star) = F(\mathbf{Q}^\star)$ とおく．
  $\mathbf{P}^\star \neq \mathbf{Q}^\star$ と仮定する．
  中点 $\mathbf{R} \defeq \tfrac{1}{2}(\mathbf{P}^\star + \mathbf{Q}^\star)$ は，
  $\CouplingsD(\mathbf{a}, \mathbf{b})$ の凸性（主張~\ref{clm:sem-coupling-convex}）よりこれに属する．
  主張~\ref{clm:sem-neg-entropy-strict-convex}（$-\Hb$ の狭義凸性）を
  中点 $\mathbf{R} = \tfrac{1}{2}\mathbf{P}^\star + \tfrac{1}{2}\mathbf{Q}^\star$ に適用すると，
  $\mathbf{P}^\star \neq \mathbf{Q}^\star$ より狭義不等号
  \[
    -\Hb(\mathbf{R})
    < \tfrac{1}{2}\bigl(-\Hb(\mathbf{P}^\star)\bigr) + \tfrac{1}{2}\bigl(-\Hb(\mathbf{Q}^\star)\bigr)
  \]
  が成り立つ．また $\inner{\mathbf{C}}{\cdot}$ は線形だから
  \[
    \inner{\mathbf{C}}{\mathbf{R}}
    = \tfrac{1}{2}\inner{\mathbf{C}}{\mathbf{P}^\star} + \tfrac{1}{2}\inner{\mathbf{C}}{\mathbf{Q}^\star}
  \]
  が等号で成り立つ．これらより（$\varepsilon > 0$ は不等号の向きを保つ）
  \begin{align*}
    F(\mathbf{R})
    &= \inner{\mathbf{C}}{\mathbf{R}} - \varepsilon\Hb(\mathbf{R}) \\
    &= \Bigl(\tfrac{1}{2}\inner{\mathbf{C}}{\mathbf{P}^\star} + \tfrac{1}{2}\inner{\mathbf{C}}{\mathbf{Q}^\star}\Bigr)
       - \varepsilon\Hb(\mathbf{R}) \\
    &< \Bigl(\tfrac{1}{2}\inner{\mathbf{C}}{\mathbf{P}^\star} + \tfrac{1}{2}\inner{\mathbf{C}}{\mathbf{Q}^\star}\Bigr)
       + \tfrac{1}{2}\bigl(-\varepsilon\Hb(\mathbf{P}^\star)\bigr)
       + \tfrac{1}{2}\bigl(-\varepsilon\Hb(\mathbf{Q}^\star)\bigr) \\
    &= \tfrac{1}{2}\bigl(\inner{\mathbf{C}}{\mathbf{P}^\star} - \varepsilon\Hb(\mathbf{P}^\star)\bigr)
       + \tfrac{1}{2}\bigl(\inner{\mathbf{C}}{\mathbf{Q}^\star} - \varepsilon\Hb(\mathbf{Q}^\star)\bigr) \\
    &= \tfrac{1}{2}F(\mathbf{P}^\star) + \tfrac{1}{2}F(\mathbf{Q}^\star) = m
  \end{align*}
  を得る．ここで第 2 行で線形性の等式を $\inner{\mathbf{C}}{\mathbf{R}}$ に代入し，
  第 3 行で $-\varepsilon\Hb(\mathbf{R})$ を狭義凸性（の $\varepsilon$ 倍）で上から押さえた．しかし $\mathbf{R} \in \CouplingsD(\mathbf{a}, \mathbf{b})$ かつ $m$ は $F$ の最小値だから
  $F(\mathbf{R}) \geq m$ でなければならず，矛盾．
  よって $\mathbf{P}^\star = \mathbf{Q}^\star$，すなわち最適解は一意である．
\end{proof}

%% ============================================================
%% 3.2 KL ダイバージェンスによる定式化
%% ============================================================
\section{KL ダイバージェンスによる定式化}
\label{sec:sem-kl-formulation}

\begin{definition}{離散 KL ダイバージェンス}{sem-kl-divergence}
  非負行列 $\mathbf{P}, \mathbf{K} \in \Rnn^{n \times m}$ に対して，
  \textbf{KL ダイバージェンス}を
  \[
    \KLD(\mathbf{P} \| \mathbf{K})
    \defeq
    \sum_{i,j}
    \left(
      P_{i,j}\log\frac{P_{i,j}}{K_{i,j}}
      - P_{i,j} + K_{i,j}
    \right)
  \]
  で定める．ただし $0\log 0 = 0$ と約束し，
  $P_{i,j} > 0$ かつ $K_{i,j}=0$ となる成分がある場合は
  $\KLD(\mathbf{P}\|\mathbf{K}) = +\infty$ とする．
\end{definition}

\begin{definition}{Gibbs カーネル}{sem-gibbs-kernel}
  コスト行列 $\mathbf{C}$ と $\varepsilon > 0$ に対して，
  \textbf{Gibbs カーネル} $\mathbf{K} \in \Rpos^{n \times m}$ を
  \[
    K_{i,j} \defeq \exp\!\left(-\frac{C_{i,j}}{\varepsilon}\right)
  \]
  で定める．
\end{definition}

\begin{proposition}{正則化 OT は KL 射影である}{sem-entropic-kl-projection}
  $\mathbf{K} = \exp(-\mathbf{C}/\varepsilon)$ とすると，
  $\MKD_{\mathbf{C}}^\varepsilon(\mathbf{a}, \mathbf{b})$ の最適解は
  \[
    \mathbf{P}_\varepsilon
    =
    \argmin_{\mathbf{P} \in \CouplingsD(\mathbf{a}, \mathbf{b})}
    \KLD(\mathbf{P}\|\mathbf{K})
  \]
  と書ける．
\end{proposition}

\begin{proof}
  $\log K_{i,j} = -C_{i,j}/\varepsilon$ より
  \begin{align*}
    \varepsilon \KLD(\mathbf{P}\|\mathbf{K})
    &=
    \varepsilon\sum_{i,j}
    \left(
      P_{i,j}\log P_{i,j}
      - P_{i,j}\log K_{i,j}
      - P_{i,j}
      + K_{i,j}
    \right) \\
    &=
    \sum_{i,j} C_{i,j}P_{i,j}
    + \varepsilon\sum_{i,j} P_{i,j}(\log P_{i,j}-1)
    + \varepsilon\sum_{i,j} K_{i,j} \\
    &=
    \inner{\mathbf{C}}{\mathbf{P}} - \varepsilon \Hb(\mathbf{P})
    + \varepsilon\sum_{i,j} K_{i,j}.
  \end{align*}
  最後の項は $\mathbf{P}$ に依存しない定数なので，
  両者の最小化問題は同じ最適解を持つ．
\end{proof}

最適解の具体的な形を導く準備として，まず最適解が領域の内部にある（全成分が正である）ことを示す．成分が $0$ にならないとわかれば，不等式制約 $\mathbf{P} \geq 0$ はどこも等号で効かないので無視でき，残る等式制約（周辺条件）だけに対して Lagrange の未定乗数法が使える．

\begin{lemma}{正則化解の正値性}{sem-entropic-positivity}
  定義~\ref{def:sem-entropic-ot} の仮定（$\mathbf{a} \in \Rpos^n$, $\mathbf{b} \in \Rpos^m$）の下で，
  命題~\ref{prop:sem-entropic-existence-unique} の一意解
  \[
    \mathbf{P}_\varepsilon
    =
    \argmin_{\mathbf{P} \in \CouplingsD(\mathbf{a}, \mathbf{b})}
    \left\{
      \inner{\mathbf{C}}{\mathbf{P}}
      - \varepsilon \Hb(\mathbf{P})
    \right\}
  \]
  は全成分が正である：
  \[
    (\mathbf{P}_\varepsilon)_{i,j} > 0 \qquad (\forall i, j).
  \]
  すなわちどの成分も $0$ に張り付かず，不等式制約 $\mathbf{P} \geq 0$ が等号で成り立つ成分は一つもない．
\end{lemma}

\begin{proof}
  命題~\ref{prop:sem-entropic-kl-projection} より，$\mathbf{P}_\varepsilon$ は
  Gibbs カーネル $\mathbf{K} = \exp(-\mathbf{C}/\varepsilon)$ への KL ダイバージェンスを
  結合集合上で最小にする点
  \[
    \mathbf{P}_\varepsilon
    =
    \argmin_{\mathbf{P} \in \CouplingsD(\mathbf{a}, \mathbf{b})} \KLD(\mathbf{P}\|\mathbf{K})
  \]
  である．用いるのは $\mathbf{K}$ の全成分が正であること（定義~\ref{def:sem-gibbs-kernel}）だけである．
  以下，$\mathbf{P}_\varepsilon$ にゼロ成分があると仮定して矛盾を導く．

  \emph{正値な比較対象．}
  積結合 $\mathbf{Q} \defeq \mathbf{a}\mathbf{b}^\top$ は，$\sum_i a_i = \sum_j b_j = 1$ より
  $\mathbf{Q}\ones_m = \mathbf{a}\,(\mathbf{b}^\top\ones_m) = \mathbf{a}$,
  $\mathbf{Q}^\top\ones_n = \mathbf{b}\,(\mathbf{a}^\top\ones_n) = \mathbf{b}$ をみたし，
  成分 $Q_{i,j} = a_i b_j > 0$ はすべて正である．
  ゆえに $\mathbf{Q} \in \CouplingsD(\mathbf{a}, \mathbf{b})$ は正値な実行可能点である．

  \emph{線分による比較．}
  $\CouplingsD(\mathbf{a}, \mathbf{b})$ の凸性（主張~\ref{clm:sem-coupling-convex}）より，$\mathbf{P}_\varepsilon$ から $\mathbf{Q}$ への線分
  \[
    \mathbf{P}^t \defeq (1-t)\mathbf{P}_\varepsilon + t\mathbf{Q}
    \in \CouplingsD(\mathbf{a}, \mathbf{b})
    \qquad (t \in [0,1]),
    \qquad
    P^t_{i,j} = (1-t)(\mathbf{P}_\varepsilon)_{i,j} + t\,Q_{i,j}
  \]
  はすべて実行可能である（始点は $\mathbf{P}^0 = \mathbf{P}_\varepsilon$）．$t \in (0,1)$ では $P^t_{i,j} \geq t\,Q_{i,j} > 0$ が全成分で正である．
  $h(t) \defeq \KLD(\mathbf{P}^t\|\mathbf{K})$ とおく．
  $\KLD = \sum_{i,j}\bigl(P_{i,j}\log\frac{P_{i,j}}{K_{i,j}} - P_{i,j} + K_{i,j}\bigr)$ の各項を成分 $P_{i,j}$ で偏微分すると
  \[
    \frac{\partial}{\partial P_{i,j}}\!\left( P_{i,j}\log\frac{P_{i,j}}{K_{i,j}} - P_{i,j} + K_{i,j} \right)
    = \log\frac{P_{i,j}}{K_{i,j}}
  \]
  であり，また線分 $P^t_{i,j} = (1-t)(\mathbf{P}_\varepsilon)_{i,j} + t\,Q_{i,j}$ より $\dfrac{d}{dt} P^t_{i,j} = Q_{i,j} - (\mathbf{P}_\varepsilon)_{i,j}$．
  よって $t \in (0,1)$（このとき $P^t_{i,j} > 0$）では，連鎖律により
  \[
    h'(t)
    = \sum_{i,j} \log\frac{P^t_{i,j}}{K_{i,j}} \,\bigl(Q_{i,j} - (\mathbf{P}_\varepsilon)_{i,j}\bigr).
  \]

  \emph{微分の発散と矛盾．}
  $h'(t)$ の和を $(\mathbf{P}_\varepsilon)_{i,j}$ の正負で二つに分け，$t \to 0+$ での挙動をみる（$K_{i,j} > 0$ は固定の正定数）．
  \begin{itemize}
    \item $(\mathbf{P}_\varepsilon)_{i,j} > 0$ の項：$P^t_{i,j} \to (\mathbf{P}_\varepsilon)_{i,j} > 0$ なので
      $\log(P^t_{i,j}/K_{i,j})$ は有限値に収束し，係数も有限．ゆえにこれらの項の和は有界である．
    \item $(\mathbf{P}_\varepsilon)_{i,j} = 0$ の項：係数 $Q_{i,j} - (\mathbf{P}_\varepsilon)_{i,j} = a_i b_j > 0$，
      かつ $\log(P^t_{i,j}/K_{i,j}) = \log(t\,a_i b_j / K_{i,j}) \to -\infty$ だから，各項は $-\infty$ に発散する．
      仮定よりこのような項は少なくとも一つある．
  \end{itemize}
  前者の和は有界，後者は $-\infty$ に発散する（$+\infty$ に飛ぶ項はないので相殺は起きない）から，
  $\lim_{t \to 0+} h'(t) = -\infty$ である．
  ゆえにある $t_1 \in (0,1)$ が存在して $(0, t_1)$ 上で $h'(t) < 0$ となり，
  $h$ は $[0, t_1]$ 上で（連続かつ）狭義減少するから，
  十分小さい $t > 0$ で $h(t) < h(0) = \KLD(\mathbf{P}_\varepsilon\|\mathbf{K})$ となる．
  これは $\mathbf{P}^t \in \CouplingsD(\mathbf{a}, \mathbf{b})$ と $\mathbf{P}_\varepsilon$ が KL 射影であることに反する．
  よって全成分は正である．
  （直観的には，ある成分を $0$ からわずかに正へ開くと，勾配 $\log(P_{i,j}/K_{i,j})$ が
  $P_{i,j} \to 0$ で $-\infty$ に落ちるため目的関数が一気に下がる．
  ゆえに $0$ 成分は最適たり得ない．）
\end{proof}

これと KL 射影としての表現（命題~\ref{prop:sem-entropic-kl-projection}）を用いると，最適解の具体的な形が定まる．

\begin{proposition}{正則化解の形（スケーリング形）}{sem-entropic-scaling-form}
  定義~\ref{def:sem-entropic-ot} の仮定（$\mathbf{a} \in \Rpos^n$, $\mathbf{b} \in \Rpos^m$）
  の下で，$\MKD_{\mathbf{C}}^\varepsilon(\mathbf{a}, \mathbf{b})$ の最適解
  $\mathbf{P}_\varepsilon$ は，ある $\mathbf{u} \in \Rpos^n$, $\mathbf{v} \in \Rpos^m$
  を用いて
  \[
    (\mathbf{P}_\varepsilon)_{i,j} = u_i\, K_{i,j}\, v_j,
    \qquad\text{すなわち}\qquad
    \mathbf{P}_\varepsilon = \diag(\mathbf{u})\,\mathbf{K}\,\diag(\mathbf{v})
  \]
  と書ける．ここで $\mathbf{K} = \exp(-\mathbf{C}/\varepsilon)$ は Gibbs カーネル
  （定義~\ref{def:sem-gibbs-kernel}）である．
  $\mathbf{P}_\varepsilon$ は一意であり，スケーリングベクトル $(\mathbf{u}, \mathbf{v})$ は，
  任意の $\lambda > 0$ による置き換え $(\mathbf{u}, \mathbf{v}) \mapsto (\lambda\mathbf{u}, \lambda^{-1}\mathbf{v})$
  の自由度を除いて一意に定まる．
\end{proposition}

\begin{proof}
  \textbf{方針．}
  補題~\ref{lem:sem-entropic-positivity} より最適解 $\mathbf{P}_\varepsilon$ は領域の内点にあり，非負制約は不活性である．
  そこで等式制約のみの 1 階条件（Lagrange の未定乗数法）を適用してスケーリング形を導き（Step 1），
  最後にスケーリングベクトルの一意性を確認する（Step 2）．

  \medskip
  \textbf{Step 1（1 階条件）．}
  補題~\ref{lem:sem-entropic-positivity} より $\mathbf{P}_\varepsilon$ は領域 $\Rpos^{n \times m}$ の内点にあるので，
  非負制約 $\mathbf{P} \geq 0$ は活性化せず，制約は等式 $\mathbf{P}\ones_m = \mathbf{a}$,
  $\mathbf{P}^\top\ones_n = \mathbf{b}$ のみと見なせる．この内点性こそが，
  以下で等式制約のみの Lagrange 条件へ移る論理的な橋渡しである．

  \emph{勾配の計算．}
  命題~\ref{prop:sem-entropic-kl-projection} より $\mathbf{P}_\varepsilon$ は
  $\KLD(\cdot \| \mathbf{K})$ の $\CouplingsD(\mathbf{a}, \mathbf{b})$ 上での最小元である．
  $\KLD(\cdot \| \mathbf{K})$ は $\Rpos^{n \times m}$ 上で微分可能で，
  定義~\ref{def:sem-kl-divergence} からの直接計算により
  \[
    \frac{\partial}{\partial P_{i,j}} \KLD(\mathbf{P} \| \mathbf{K})
    = \log\frac{P_{i,j}}{K_{i,j}}
  \]
  をみたす．

  \emph{Lagrange 1 階条件．}
  制約はすべてアフィン（一次）だから制約想定（constraint qualification）がみたされ，
  内点の最小元では Lagrange の 1 階条件が必要条件となる．
  ここで，行和制約 $(\mathbf{P}\ones_m)_i = a_i$ の $P_{k,l}$ に関する偏微分は $[k = i]$，
  列和制約 $(\mathbf{P}^\top\ones_n)_j = b_j$ の偏微分は $[l = j]$ である
  （$[\,\cdot\,]$ は指示関数）．
  すなわち行制約は $i$ 行にのみ $1$ を，列制約は $j$ 列にのみ $1$ を立てる．
  ゆえに乗数 $\alpha \in \R^n$, $\beta \in \R^m$（符号制約はない）が存在して，
  停留条件は
  \[
    \log\frac{(\mathbf{P}_\varepsilon)_{i,j}}{K_{i,j}}
    = \alpha_i \cdot 1 + \beta_j \cdot 1 = \alpha_i + \beta_j
    \qquad (\forall i, j)
  \]
  という行成分 $\alpha_i$ と列成分 $\beta_j$ の和の形をとる．

  \emph{指数化．}
  両辺の指数をとり $u_i \defeq e^{\alpha_i} > 0$, $v_j \defeq e^{\beta_j} > 0$ とおけば
  （正値性は指数関数から自動的に従う），
  \[
    (\mathbf{P}_\varepsilon)_{i,j} = u_i K_{i,j} v_j,
    \qquad\text{すなわち}\qquad
    \mathbf{P}_\varepsilon = \diag(\mathbf{u})\,\mathbf{K}\,\diag(\mathbf{v})
  \]
  を得る．これが命題本文のスケーリング形である．

  \medskip
  \textbf{Step 2（一意性）．}
  $\mathbf{P}_\varepsilon$ 自体の一意性は命題~\ref{prop:sem-entropic-existence-unique}
  （正則化解の存在と一意性）による（ここでは再証明しない）．
  スケーリングベクトル $(\mathbf{u}, \mathbf{v})$ については，次のとおり $\lambda$ 自由度を除いて一意である．

  まず十分性（自由度の存在）として，任意の $\lambda > 0$ に対し
  $(\lambda u_i)(\lambda^{-1} v_j) = u_i v_j$ より $(\lambda\mathbf{u}, \lambda^{-1}\mathbf{v})$ も
  同じ $\mathbf{P}_\varepsilon$ を与える．
  逆にこの自由度を除けば一意であることを，補助記号
  $w_{i,j} \defeq u_i v_j$ を用いて 3 段で示す．
  \begin{itemize}
    \item[(i)] 全成分が正なので，各 $(i,j)$ で積
      $w_{i,j} = u_i v_j = (\mathbf{P}_\varepsilon)_{i,j} / K_{i,j} > 0$ が一意に定まる
      （個々の $u_i, v_j$ は不定だが，積 $w_{i,j}$ は確定する）．
    \item[(ii)] 比 $u_i / u_{i'} = w_{i,j} / w_{i',j}$ は $j$ に依らない
      （分子分母で $v_j$ が約分されて消える；
      これは $\log w_{i,j} = \alpha_i + \beta_j$ で $\alpha_i$ が $j$ に依らないことに対応する）．
    \item[(iii)] ゆえに $\mathbf{u}$ は正の定数倍を除いて一意であり，
      対応して $v_j = w_{i,j} / u_i$ も定まる．
  \end{itemize}
\end{proof}

\begin{remark}{正則化解の指数表示}{sem-entropic-exp-form}
  Step 1 の停留条件は，乗数を $f_i \defeq \varepsilon\alpha_i$, $g_j \defeq \varepsilon\beta_j$ と
  スケールし直すと，元のコスト $\mathbf{C}$ を用いた指数表示
  \[
    (\mathbf{P}_\varepsilon)_{i,j}
    = \exp\!\left(\frac{f_i + g_j - C_{i,j}}{\varepsilon}\right)
    = e^{f_i / \varepsilon}\, K_{i,j}\, e^{g_j / \varepsilon}
  \]
  に書き換えられる（$u_i = e^{f_i / \varepsilon}$, $v_j = e^{g_j / \varepsilon}$）．
  すなわち「コストとポテンシャルの差 $f_i + g_j - C_{i,j}$ を $\varepsilon$ で割って指数化する」形である．
  この式は，KL 射影（命題~\ref{prop:sem-entropic-kl-projection}）を経由せず，
  元の目的関数 $F(\mathbf{P}) = \inner{\mathbf{C}}{\mathbf{P}} - \varepsilon\Hb(\mathbf{P})$ に
  直接 1 階条件（$\partial(-\Hb)/\partial P_{i,j} = \log P_{i,j}$，注意~\ref{rem:sem-entropy-convention}）
  を当てても同じく得られ，スケーリング形が定式化の選び方に依らないことを示す．
  なお $\mathbf{f}, \mathbf{g}$ は周辺制約に付随する双対変数の役割をもつが，双対理論は後の章で扱う．
\end{remark}

\begin{remark}{連続版：Schr\"odinger 問題}{sem-entropic-schrodinger}
  連続版（$\X, \Y$ を Polish 空間とする）でも対応する主張が成り立ち，最適カップリング $\pi$ は
  基準測度に対する密度が $u(x)\, e^{-c(x,y)/\varepsilon}\, v(y)$ という積の形をとる．
  この問題は \textbf{Schr\"odinger 問題}（1931）と呼ばれ，離散の 1 階条件
  $\log\bigl((\mathbf{P}_\varepsilon)_{i,j} / K_{i,j}\bigr) = \alpha_i + \beta_j$ に対応する停留条件は，
  変分法（オイラー–ラグランジュ方程式）によって与えられる．
  本セミナーでは離散版に限って扱う．
\end{remark}

%% ============================================================
%% 3.3 正則化パラメータの極限
%% ============================================================
\section{正則化パラメータの極限}
\label{sec:sem-epsilon-limits}

非正則化 Kantorovich 問題が複数の最適解を持つとき，
エントロピー正則化はそのうちの一つ——
エントロピーが最大のもの——を自然に選び出す．
正則化パラメータ $\varepsilon$ の二つの極限はそれぞれ異なる解を選ぶ：
$\varepsilon \to 0$ では非正則化最適解の中で最大エントロピーをもつものに，
$\varepsilon \to +\infty$ では周辺分布のみで定まる独立カップリングに収束する．

\begin{claim}{独立カップリングはエントロピー最大解}{sem-max-entropy-independent}
  独立カップリング $\mathbf{a}\mathbf{b}^\top = (a_i b_j)_{i,j}$ は
  $\CouplingsD(\mathbf{a},\mathbf{b})$ に属し，その上で離散エントロピー $\Hb$ を
  最大化する唯一の元である：
  \[
    \mathbf{a}\mathbf{b}^\top
    = \argmax_{\mathbf{P} \in \CouplingsD(\mathbf{a}, \mathbf{b})} \Hb(\mathbf{P}).
  \]
\end{claim}

\begin{proof}
  \textbf{$\mathbf{a}\mathbf{b}^\top \in \CouplingsD(\mathbf{a},\mathbf{b})$ であること．}
  $\sum_i a_i = \sum_j b_j = 1$ より
  $(\mathbf{a}\mathbf{b}^\top)\ones_m = \mathbf{a}(\mathbf{b}^\top\ones_m) = \mathbf{a}$，
  $(\mathbf{a}\mathbf{b}^\top)^\top\ones_n = \mathbf{b}(\mathbf{a}^\top\ones_n) = \mathbf{b}$．
  また $\mathbf{a}, \mathbf{b} \in \Rpos$ より $a_i b_j > 0$．

  \textbf{KL ダイバージェンスによる表示．}
  $\mathbf{K} \defeq \mathbf{a}\mathbf{b}^\top \in \Rpos^{n \times m}$ とおく．
  任意の $\mathbf{P} \in \CouplingsD(\mathbf{a},\mathbf{b})$ に対し，
  周辺条件 $\sum_j P_{i,j} = a_i$, $\sum_i P_{i,j} = b_j$ と
  $\sum_{i,j} P_{i,j} = \sum_{i,j} a_i b_j = 1$ を用いると
  \begin{align*}
    \KLD(\mathbf{P}\|\mathbf{K})
    &= \sum_{i,j}
       \left( P_{i,j}\log\frac{P_{i,j}}{a_i b_j} - P_{i,j} + a_i b_j \right) \\
    &= \sum_{i,j} P_{i,j}\log P_{i,j}
       - \sum_i a_i \log a_i - \sum_j b_j \log b_j \\
    &= \bigl(1 - \Hb(\mathbf{P})\bigr)
       - \sum_i a_i \log a_i - \sum_j b_j \log b_j
  \end{align*}
  となる．第 2 等号は
  $-\sum_{i,j}P_{i,j} + \sum_{i,j}a_i b_j = -1 + 1 = 0$ と，周辺条件による
  $\sum_{i,j} P_{i,j}\log(a_i b_j) = \sum_i a_i\log a_i + \sum_j b_j\log b_j$ から従う．
  第 3 等号は
  $\Hb(\mathbf{P}) = -\sum_{i,j}P_{i,j}(\log P_{i,j}-1) = 1 - \sum_{i,j}P_{i,j}\log P_{i,j}$
  （$\sum_{i,j}P_{i,j}=1$）による．したがって
  \[
    \Hb(\mathbf{P})
    = \underbrace{\,1 - \sum_i a_i \log a_i - \sum_j b_j \log b_j\,}
        _{\mathbf{P} \text{ に依らない定数}}
      - \KLD(\mathbf{P}\|\mathbf{K}).
  \]

  \textbf{非負性と一意性．}
  各成分について $g(p) \defeq p\log(p/k) - p + k$（$k > 0$）は
  $g'(p) = \log(p/k)$ より $p = k$ で最小値 $g(k) = 0$ をとり，
  $g \geq 0$ である（$g(0) = k > 0$）．
  ゆえに $\KLD(\mathbf{P}\|\mathbf{K}) \geq 0$ であり，
  等号は全成分で $P_{i,j} = a_i b_j$，すなわち $\mathbf{P} = \mathbf{a}\mathbf{b}^\top$ のときに限る．
  上の表示より $\Hb$ は $\mathbf{P} = \mathbf{a}\mathbf{b}^\top$ で唯一の最大値をとる．
\end{proof}

\begin{theorem}{正則化パラメータ $\varepsilon$ の極限による収束}{sem-entropic-convergence-eps-zero}
  各 $\varepsilon > 0$ に対し，$\mathbf{P}_\varepsilon$ を
  $\MKD_{\mathbf{C}}^\varepsilon(\mathbf{a}, \mathbf{b})$ の一意解とする．
  このとき，以下が成り立つ．

  \begin{enumerate}
    \item[(i)]
      最適値は収束する：
      $\displaystyle
        \lim_{\varepsilon \to 0}
        \MKD_{\mathbf{C}}^\varepsilon(\mathbf{a},\mathbf{b})
        =
        \MKD_{\mathbf{C}}(\mathbf{a},\mathbf{b}).
      $
    \item[(ii)]
      $\mathbf{P}_\varepsilon$ の $\varepsilon \to 0$ における任意の極限点は
      非正則化問題の最適解である．
    \item[(iii)]
      $\{\mathbf{P}_\varepsilon\}_{\varepsilon > 0}$ 全体が，
      非正則化最適解の中でエントロピー $\Hb$ を最大化する
      一意な解に収束する：
      \[
        \mathbf{P}_\varepsilon
        \;\xrightarrow{\varepsilon \to 0}\;
        \argmax_{\mathbf{P} \in S^*} \Hb(\mathbf{P}).
      \]
    \item[(iv)]
      $\varepsilon \to +\infty$ のとき，$\mathbf{P}_\varepsilon$ は
      独立カップリングに収束する：
      \[
        \mathbf{P}_\varepsilon
        \;\xrightarrow{\varepsilon \to +\infty}\;
        \mathbf{a}\mathbf{b}^\top = (a_i b_j)_{i,j}.
      \]
  \end{enumerate}
\end{theorem}

\begin{proof}
  非正則化問題の任意の最適解を $\mathbf{P}^*$ とする．

  \textbf{(i), (ii).}
  $\mathbf{P}_\varepsilon$ の最適性より
  \[
    \inner{\mathbf{C}}{\mathbf{P}_\varepsilon}
    - \varepsilon\Hb(\mathbf{P}_\varepsilon)
    \leq
    \inner{\mathbf{C}}{\mathbf{P}^*}
    - \varepsilon\Hb(\mathbf{P}^*).
  \]
  一方，$\mathbf{P}^*$ は非正則化問題の最適解なので
  $\inner{\mathbf{C}}{\mathbf{P}^*}
  \leq \inner{\mathbf{C}}{\mathbf{P}_\varepsilon}$ であり，
  \[
    0
    \leq
    \inner{\mathbf{C}}{\mathbf{P}_\varepsilon}
    - \inner{\mathbf{C}}{\mathbf{P}^*}
    \leq
    \varepsilon\bigl(\Hb(\mathbf{P}_\varepsilon)-\Hb(\mathbf{P}^*)\bigr).
  \]
  $\CouplingsD(\mathbf{a}, \mathbf{b})$ はコンパクトであり，$\Hb$ はその上で有界なので，
  右辺は $\varepsilon \to 0$ で $0$ に収束する．
  したがって
  $\inner{\mathbf{C}}{\mathbf{P}_\varepsilon}
  \to \inner{\mathbf{C}}{\mathbf{P}^*}
  = \MKD_{\mathbf{C}}(\mathbf{a},\mathbf{b})$ であり，
  $\varepsilon\Hb(\mathbf{P}_\varepsilon) \to 0$ なので
  $\MKD_{\mathbf{C}}^\varepsilon \to \MKD_{\mathbf{C}}$ を得る．
  コンパクト集合 $\CouplingsD(\mathbf{a},\mathbf{b})$ 内の
  任意の極限点は非正則化問題の最適値を達成するから，
  最適解である．

  \textbf{(iii).}
  上の不等式から
  $\Hb(\mathbf{P}_\varepsilon) \geq \Hb(\mathbf{P}^*)$
  が任意の非正則化最適解 $\mathbf{P}^*$ について成り立つ．
  $\Hb$ は連続であるから，任意の極限点 $\mathbf{P}^0$ は
  非正則化最適解の中で $\Hb$ を最大化する．
  最適解集合 $S^*$ は凸であり（主張~\ref{clm:sem-optimal-face}），
  主張~\ref{clm:sem-neg-entropy-strict-convex} より $-\Hb$ は
  （凸部分集合 $S^*$ 上でも）狭義凸であるから，
  命題~\ref{prop:sem-strict-convex-unique-min} より
  $\Hb$ を最大化する $S^*$ の元は一意である．
  したがって $\{\mathbf{P}_\varepsilon\}$ の任意の収束部分列は
  同一の極限を持ち，コンパクト集合内の点列でこの性質を持つものは
  全体として収束するから，
  $\mathbf{P}_\varepsilon$ は最大エントロピー最適解に収束する．

  \textbf{(iv).}
  $\varepsilon > 0$ で目的関数を割れば，$\mathbf{P}_\varepsilon$ は
  \[
    \mathbf{P}_\varepsilon
    = \argmin_{\mathbf{P} \in \CouplingsD(\mathbf{a},\mathbf{b})}
      \left( \tfrac{1}{\varepsilon}\inner{\mathbf{C}}{\mathbf{P}} - \Hb(\mathbf{P}) \right)
  \]
  の最小元でもある．$\varepsilon_\ell \to +\infty$ なる任意の列をとり，
  $\CouplingsD(\mathbf{a},\mathbf{b})$ のコンパクト性より収束部分列
  $\mathbf{P}_{\varepsilon_\ell} \to \mathbf{P}^\infty \in \CouplingsD(\mathbf{a},\mathbf{b})$ を抽出する．
  任意の $\mathbf{Q} \in \CouplingsD(\mathbf{a},\mathbf{b})$ について最小性から
  \[
    \tfrac{1}{\varepsilon_\ell}\inner{\mathbf{C}}{\mathbf{P}_{\varepsilon_\ell}} - \Hb(\mathbf{P}_{\varepsilon_\ell})
    \leq
    \tfrac{1}{\varepsilon_\ell}\inner{\mathbf{C}}{\mathbf{Q}} - \Hb(\mathbf{Q})
  \]
  が成り立つ．$\inner{\mathbf{C}}{\cdot}$ はコンパクト集合上で有界だから
  $\varepsilon_\ell^{-1}\inner{\mathbf{C}}{\cdot} \to 0$ であり，$\Hb$ は連続なので，
  $\ell \to \infty$ の極限で $\Hb(\mathbf{P}^\infty) \geq \Hb(\mathbf{Q})$ が
  すべての $\mathbf{Q} \in \CouplingsD(\mathbf{a},\mathbf{b})$ について成り立つ．
  ゆえに $\mathbf{P}^\infty$ は $\CouplingsD(\mathbf{a},\mathbf{b})$ 上で $\Hb$ を最大化し，
  主張~\ref{clm:sem-max-entropy-independent} より $\mathbf{P}^\infty = \mathbf{a}\mathbf{b}^\top$．
  極限は部分列のとり方によらず一意だから，(iii) と同様に
  コンパクト集合内での全体収束により
  $\mathbf{P}_\varepsilon \to \mathbf{a}\mathbf{b}^\top$（$\varepsilon \to +\infty$）を得る．
\end{proof}
